\subsection{Associativity}

The following lemma completes Prop. 11 of Ch. V, §8, No.~5: \(^{2}\)

Lemma 1. --- For \(1 \leqslant i \leqslant n\) , let \(\mathrm{X}_i, \mathrm{Y}_i\) be two locally compact spaces, \(\mu_i\) a measure on \(\mathrm{X}_i\) , and \(\varphi_i\) a continuous mapping of \(\mathrm{X}_i\) into \(\mathrm{Y}_i\) . Let \(\mathrm{X} = \prod_i \mathrm{X}_i\) , \(\mathrm{Y} = \prod_i \mathrm{Y}_i\) , \(\mu = \bigotimes_i \mu_i\) , and \(\varphi\) the mapping of \(\mathrm{X}\) into \(\mathrm{Y}\) that is the product of the \(\varphi_i\) . If \(\varphi\) is \(\mu\) -proper and \(\mu_i \neq 0\) for each \(i\) , then the \(\varphi_i\) are \(\mu_i\) -proper and \(\varphi(\mu) = \bigotimes_i \varphi_i(\mu_i)\) .

We can suppose that the \(\mu_{i}\) are positive and \(n = 2\) . Let \(f_{1} \in \mathcal{K}_{+}(Y_{1})\) . Since \(\mu_{2} \neq 0\) , there exists an \(f_{2} \in \mathcal{K}_{+}(Y_{2})\) such that \(f_{2} \circ \varphi_{2}\) is not \(\mu_{2}\) -negligible. The function \((x_{1}, x_{2}) \mapsto f_{1}(\varphi_{1}(x_{1})) f_{2}(\varphi_{2}(x_{2}))\) is essentially \(\mu\) -integrable and continuous, hence \(\mu\) -integrable. Therefore there exists an \(x_{2} \in X_{2}\) such that \(f_{2}(\varphi_{2}(x_{2})) \neq 0\) and such that the function \(x_{1} \mapsto f_{1}(\varphi_{1}(x_{1})) f_{2}(\varphi_{2}(x_{2}))\) is \(\mu_{1}\) -integrable. Therefore \(f_{1} \circ \varphi_{1}\) is \(\mu_{1}\) -integrable, which proves that \(\varphi_{1}\) is \(\mu_{1}\) -proper. One argues similarly for \(\varphi_{2}\) . Then \(\varphi(\mu) = \bigotimes_{i} \varphi_{i}(\mu_{i})\) by Prop. 11 of Ch. V, §8, No.~5.

The following lemma completes Prop. 4 of Ch. V, §6, No.~3. \(^{3}\)

Lemma 2. --- Let \(\mathrm{T}, \mathrm{T}', \mathrm{T}''\) be three locally compact spaces, \(\mu\) a measure on \(\mathrm{T}\) , \(\pi\) a \(\mu\) -measurable mapping of \(\mathrm{T}\) into \(\mathrm{T}'\) , \(\pi'\) a continuous mapping of \(\mathrm{T}'\) into \(\mathrm{T}''\) , and \(\pi'' = \pi' \circ \pi\) . If \(\pi''\) is \(\mu\) -proper, then \(\pi\) is \(\mu\) -proper, \(\pi'\) is \(\pi(\mu)\) -proper, and \(\pi''(\mu) = \pi' (\pi(\mu))\) .

Let \(K'\) be a compact subset of \(T'\) . Then \(K'' = \pi'(K')\) is compact, therefore \(\pi''(K'')\) is essentially \(\mu\) -integrable, therefore \(\overline{\pi}(K') \subset \overline{\pi''}(K'')\) is essentially \(\mu\) -integrable, thus \(\pi\) is \(\mu\) -proper. Then \(\pi'\) is \(\pi(\mu)\) -proper and \(\pi''(\mu) = \pi'(\pi(\mu))\) by Ch. V, §6, No.~3, Prop. 4.

PROPOSITION 1. --- Let \(\mathrm{X}_{ij}\) ( \(1 \leqslant i \leqslant m, 1 \leqslant j \leqslant n_i\) ), \(\mathrm{Y}_i\) ( \(1 \leqslant i \leqslant m\) ), and \(\mathbf{Z}\) be locally compact spaces; for each \(i\) , let \(\varphi_i\) be a mapping of \(\mathrm{X}_i = \prod_j \mathrm{X}_{ij}\) into \(\mathrm{Y}_i\) ; let \(\varphi\) be the product of the \(\varphi_i\) , mapping \(\mathrm{X} = \prod_i \mathrm{X}_i\) into \(\mathrm{Y} = \prod_j \mathrm{Y}_i\) ; let \(\psi\) be a mapping of \(\mathrm{Y}\) into \(\mathbf{Z}\) .

\begin{enumerate}
\def\labelenumi{(\roman{enumi})}
\tightlist
\item
  Let \(\mu_{ij}\) be measures given, respectively, on the \(\mathrm{X}_{ij}\) , such that, for each \(i\) , the \(\mu_{ij}\) ( \(1 \leqslant j \leqslant n_i\) ) are \(\varphi_i\) -convolvable, and such that the measures \(*|_{\mu_{ij}}|\) are \(\psi\) -convolvable; then the \(\mu_{ij}\) , for \(1 \leqslant i \leqslant m\) , \(1 \leqslant j \leqslant n_i\) , are \((\psi \circ \varphi)\) -convolvable and
\end{enumerate}

\[
\underset {i, j} {\ast} \mu_ {i j} = \underset {i} {\ast} \left(\underset {j} {\ast} \mu_ {i j}\right).\tag{3}
\]

\begin{enumerate}
\def\labelenumi{(\roman{enumi})}
\setcounter{enumi}{1}
\tightlist
\item
  Assume \(\psi\) and the \(\varphi_{i}\) continuous, and let \(\mu_{ij}\) be measures \(\neq 0\) given, respectively, on the \(X_{ij}\) and \((\psi \circ \varphi)\) -convolvable; then, for each \(i\) , the \(\mu_{ij}\) ( \(1 \leqslant j \leqslant n_{i}\) ) are \(\varphi_{i}\) -convolvable, the measures \(*_{j} |\mu_{ij}|\) are \(\psi\) -convolvable, and the formula (3) holds.
\end{enumerate}

It suffices to consider the case that all of the measures in question are \(\geqslant 0\) .

Let us place ourselves under the hypotheses of (i). The mapping \(\varphi\) is proper for \(\bigotimes_{i,j} \mu_{ij}\) , and

\[
\varphi \left(\bigotimes_ {i, j} \mu_ {i j}\right) = \bigotimes_ {i} \varphi_ {i} \left(\bigotimes_ {j} \mu_ {i j}\right) = \bigotimes_ {i} \left(* _ {j} \mu_ {i j}\right)
\]

(Ch. V, §8, No.~5, Prop. 11). The mapping \(\psi \circ \varphi\) is proper for \(\bigotimes_{i,j} \mu_{ij}\) , and

\[
(\psi \circ \varphi) \left(\bigotimes_ {i, j} \mu_ {i j}\right) = \psi \left(\bigotimes_ {i} \left(* _ {j} \mu_ {i j}\right)\right) = * _ {i} \left(* _ {j} \mu_ {i j}\right)
\]

(Ch. V, §6, Prop. 4). Therefore the \(\mu_{ij}\) ( \(1 \leqslant i \leqslant m\) , \(1 \leqslant j \leqslant n_i\) ) are \((\psi \circ \varphi)\) -convolvable and formula (3) holds.

Let us place ourselves under the hypotheses of (ii). First of all, Lemma 2 proves that \(\varphi\) is proper for \(\bigotimes_{i,j} \mu_{ij}\) . Lemma 1 then proves that for every \(i\) , \(\varphi_i\) is proper for \(\bigotimes_{i} \mu_{ij}\) , and that

\[
\varphi \Big (\bigotimes_ {i, j} \mu_ {i j} \Big) = \bigotimes_ {i} \Big (* _ {j} \mu_ {i j} \Big).
\]

By Lemma 2, \(\psi\) is proper for \(\bigotimes_{i}\left(*_{j}\mu_{ij}\right)\) . Whence the proposition.

COROLLARY. --- Let \(X_{i}, X_{i}^{\prime}\) ( \(1 \leqslant i \leqslant n\) ), Y, \(Y^{\prime}\) be locally compact spaces; let \(\varphi, \varphi^{\prime}\) be continuous mappings of \(X = \prod_{i} X_{i}\) into Y and of \(X^{\prime} = \prod_{i} X_{i}^{\prime}\) into \(Y^{\prime}\) , respectively; let \(f_{i}\) be continuous mappings of \(X_{i}\) into \(X_{i}^{\prime}\) ( \(1 \leqslant i \leqslant n\) ) and g a continuous mapping of Y into \(Y^{\prime}\) , such that \(\varphi^{\prime} \circ f = g \circ \varphi\) , f being the mapping of X into \(X^{\prime}\) that is the product of the \(f_{i}\) . Let \(\mu_{i}\) be measures given respectively on the \(X_{i}\) , all \(\neq 0\) . Then the following two assertions are equivalent:

\begin{enumerate}
\def\labelenumi{(\roman{enumi})}
\item
  \(f_{i}\) is \(\mu_{i}\) -proper for all \(i\) , and the measures \(f_{i}(|\mu_{i}|)\) are \(\varphi'\) -convolvable;
\item
  the \(\mu_{i}\) are \(\varphi\) -convolvable, and \(g\) is proper for \(*_{\varphi}(|\mu_i|)\) .
\end{enumerate}

Moreover, when these assertions are verified,

\[
\left. \ast_ {\varphi^ {\prime}} \left(f _ {i} \left(\mu_ {i}\right)\right) = g \left(\ast_ {\varphi} \left(\mu_ {i}\right)\right) = \ast_ {g \circ \varphi} \left(\mu_ {i}\right). \right.\tag{4}
\]

For, let \(h = \varphi' \circ f = g \circ \varphi\) . By Prop. 1, the conditions (i) and (ii) are each equivalent to the following condition:

\begin{enumerate}
\def\labelenumi{(\roman{enumi})}
\setcounter{enumi}{2}
\tightlist
\item
  the \(\mu_{i}\) are \(h\) -convolvable.
\end{enumerate}

When this is so,

\[
\left. \ast_ {\varphi^ {\prime}} \left(f _ {i} \left(\mu_ {i}\right)\right) = \ast_ {h} \left(\mu_ {i}\right) = g \left(\ast_ {\varphi} \left(\mu_ {i}\right)\right). \right.
\]

